% Generated by scripts/assemble.py; corrections belong in source/corrections.json.
\examhead{76}{吉林大学 2025 年数学分析试卷}{74}
% p075-b22; source PDF page 75
\par\Needspace{5\baselineskip}
\blockanchor{p075-b22}
\majorquestion{一}计算题.\par
% p075-b23; source PDF page 75
\blockanchor{p075-b23}
\question{1}$\lim_{x\to \pi}\frac{\sin(mx)}{\sin((2n + 1)x)}$ ，其中 $m,n$ 为正整数\par
% p075-b24; source PDF page 75
\blockanchor{p075-b24}
\question{2}$\lim_{n\to\infty}\left(\frac{2}{\pi}\arctan n\right)^{n}.$\par
% p076-b00; source PDF page 76
\blockanchor{p076-b00}
\question{3}$\lim_{x\to 0}\frac{\sqrt{\frac{1}{4} - x} + \sqrt[3]{\frac{1}{8} + x^2} - e^{2x}}{\sin x}.$\par
% p076-b01; source PDF page 76
\blockanchor{p076-b01}
\question{4}$\lim_{n\to\infty}\left(\frac{n}{n^{2}+1}+\frac{n}{n^{2}+2^{2}}+\cdots+\frac{n}{n^{2}+n^{2}}\right).$\par
% p076-b02; source PDF page 76
\blockanchor{p076-b02}
\question{5}$\lim_{n\to \infty}n\int_0^1 x^n\arctan x\mathrm{d}x.$\par
% p076-b03; source PDF page 76
\blockanchor{p076-b03}
\question{6}设 $u = \frac{x + z}{y + z}$，其中 $z$ 为由方程 $ze^z = xe^x + ye^y$ 所确定的 $x, y$ 的函数，求 $\frac{\partial u}{\partial x}, \frac{\partial u}{\partial y}$.\par
% p076-b04; source PDF page 76
\blockanchor{p076-b04}
\question{7}计算重积分 $\iiint_{D} \sqrt{x^2 + y^2} \, \mathrm{d}x \, \mathrm{d}y \, \mathrm{d}z$, 其中 $D$ 是 $x^2 + y^2 = z^2$ 与 $z = 1$ 所围的区域.\par
% p076-b05; source PDF page 76
\par\Needspace{6\baselineskip}
\blockanchor{p076-b05}
\question{8}计算第二型曲面积分\par
\nopagebreak[4]
% p076-b06; source PDF page 76
\blockanchor{p076-b06}
\displayformula{I = \iint_ {\Sigma} x y ^ {3} \mathrm{d} z \mathrm{d} z + 2 x ^ {3} y \mathrm{d} z \mathrm{d} x + 3 z ^ {2} \mathrm{d} x \mathrm{d} y.}
% p076-b07; source PDF page 76
\blockanchor{p076-b07}
其中曲面 $\Sigma$ 是锥面 $z^2 = x^2 + y^2$ 在 $z = 0$ 和 $z = 2$ 之间的部分, 方向取外侧.\par
% p076-b08; source PDF page 76
\par\Needspace{5\baselineskip}
\blockanchor{p076-b08}
\majorquestion{二}求级数 $\sum_{n=1}^{\infty}(n^{2}-1)x^{n}$ 的和函数.\par
% p076-b09; source PDF page 76
\par\Needspace{5\baselineskip}
\blockanchor{p076-b09}
\majorquestion{三}讨论当 $p > 0$ 为何值时, 级数 $\sum_{n=2}^{\infty} \ln \left(1 + \frac{(-1)^n}{n^p}\right)$ 绝对收敛, 条件收敛, 发散.\par
% p076-b10; source PDF page 76
\par\Needspace{5\baselineskip}
\blockanchor{p076-b10}
\majorquestion{四}设函数 $f(x)$ 在 $[0, +\infty)$ 上连续，且 $\lim_{x \to +\infty} f(x)$ 存在. 证明：$f(x)$ 在 $[0, +\infty)$ 上有界，并且最大值和最小值至少有一个能取到.\par
% p076-b11; source PDF page 76
\par\Needspace{5\baselineskip}
\blockanchor{p076-b11}
\majorquestion{五}设函数 $f(x)$ 在 $[0,1]$ 上连续, 在 $(0,1)$ 内可导, $f(0)=0,f(1)=1$. 证明:\par
% p076-b12; source PDF page 76
\par\Needspace{6\baselineskip}
\blockanchor{p076-b12}
存在 $\xi,\eta\in(0,1),\xi\ne\eta$, 使得\par
\nopagebreak[4]
% p076-b13; source PDF page 76
\blockanchor{p076-b13}
\displayformula{f ^ {\prime} (\xi) f ^ {\prime} (\eta) = 1.}
% p076-b14; source PDF page 76
\par\Needspace{5\baselineskip}
\blockanchor{p076-b14}
\majorquestion{六}设函数 $f(\alpha, \beta) = \int_{\Gamma} \frac{x \, \mathrm{d}y - y \, \mathrm{d}x}{\alpha x^2 + \beta y^2}, \alpha, \beta > 0$, 其中 $\Gamma$ 为圆周 $x^2 + y^2 = 1$, 取逆时针方向.\par
% p076-b15; source PDF page 76
\blockanchor{p076-b15}
(1) 求 $f(\alpha, \beta)$ 的表达式.\par
% p076-b16; source PDF page 76
\par\Needspace{6\baselineskip}
\blockanchor{p076-b16}
(2) 讨论广义重积分\par
\nopagebreak[4]
% p076-b17; source PDF page 76
\blockanchor{p076-b17}
\displayformula{\iint_ {D} f (\alpha , \beta) ^ {p} \mathrm{d} \alpha \mathrm{d} \beta}
% p076-b18; source PDF page 76
\blockanchor{p076-b18}
的敛散性, 其中 $p \in (-\infty, +\infty)$, 区域 $D = \{(\alpha, \beta) \in \mathbb{R}^2 \mid \alpha^2 + \beta^2 \leq 1, \alpha, \beta \geq 0\}$.\par
% p076-b19; source PDF page 76
\par\Needspace{5\baselineskip}
\blockanchor{p076-b19}
\majorquestion{七}证明:\par
% p076-b20; source PDF page 76
\blockanchor{p076-b20}
(1) $f(x)$ 在 $(a, b)$ 上一致连续当且仅当 $f(x)$ 在 $(a, b)$ 上连续，并且极限 $\lim_{x \to a^{+}} f(x)$ 和 $\lim_{x \to b^{-}} f(x)$ 都存在.\par
% p076-b21; source PDF page 76
\blockanchor{p076-b21}
(2) 如果 $f(x)$ 在 $(0,1]$ 上可导, 并且极限 $\lim_{x\to 0^{+}}x^{2 / 3}f'(x)$ 存在, 则 $f(x) + \frac{\sin x}{x}$ 在 $(0,1]$ 上一致连续.\par
